\section{General comments}
blablabla



\section*{Official Review by Reviewer QbaF}

\textbf{Summary.}

This paper studies online learning to defer under partial feedback and non-stationary conditions, where only the residual of the selected expert is observed at each round. It models expert residuals using a factorized switching linear dynamical system with regime changes, a shared global factor, and expert-specific states. The shared factor links experts, so enables cross-expert information transfer. The routing decision leverages predictive residual distributions with an information-directed sampling to balance exploration and exploitation.

\textbf{Strengths.}
\begin{enumerate}
    \item Clear structured modeling of expert residual and good presentation of this model.
    \item The authors deal with non-stationarity by leveraging factorized switching linear dynamical system which is also enabling closed-form updates and making the method computationally feasible.
    \item Results on synthetic and real-world datasets are consistent with the paper's claims.
    \item Figure 2 is a good presentation for understanding this model.
\end{enumerate}

\textbf{Weaknesses.}
\begin{enumerate}
    \item While the proposed model structure is interesting, the paper provides limited discussion on statistical guarantees. So it is difficult to compare the proposed method with existing approaches.
    \item The computational cost of the proposed method is not clearly analyzed.
    \item The experimental comparison is limited to LinUCB and NeuralUCB. It is unclear why the method is not compared with other exploration strategies such as Thompson sampling or ensemble sampling (Lu and Van Roy, 2017).
\end{enumerate}

Lu, Xiuyuan, and Benjamin Van Roy. ``Ensemble sampling.'' \textit{Advances in neural information processing systems} 30 (2017).

\textbf{Key Questions For Author.}
\begin{enumerate}
    \item Can you provide a clear discussion of the computational complexity of the proposed method? Also, it would be helpful to understand how the runtime compares to the baselines used in the experiments.
    \item In the LinUCB setting, the paper uses a separate parameter for each expert. It is unclear why a shared parameter is not considered. The standard contextual LinUCB typically assumes a common parameter across arms(experts). It would be helpful to understand the motivation for this choice, and whether comparisons with a shared-parameter LinUCB would change the results.
\end{enumerate}

\textbf{Limitations:} yes

\textbf{Rating:} 4: Weak accept: Technically solid paper that advances at least one sub-area of AI, with a contribution that others are likely to build on, but with some weaknesses that limit its impact (e.g., limited evaluation). Please use sparingly.

\textbf{Confidence:} 1: Your assessment is an educated guess. The submission is not in your area, or the submission was difficult to understand. Math/other details were not carefully checked.

\subsection{Response}

We thank the reviewer for the positive assessment of the model structure, the tractable filtering design, and the overall empirical direction. We respond point by point below.

\begin{itemize}
    \item \textbf{Weakness 1: limited discussion on statistical guarantees.}
    \begin{itemize}
        \item We agree that the current paper does not provide regret-style or minimax guarantees, and that this is a limitation of scope. Our goal in this work was to study the full routing problem with \emph{simultaneously} non-stationary latent regimes, partial feedback, heterogeneous experts, and time-varying expert availability. In our view, a meaningful regret analysis in this full setting is substantially more delicate than in standard stationary contextual bandits. A simplified regret theorem would likely require removing one or more of these ingredients, thereby materially changing the problem studied here. For this reason, we position the submission as a methodology paper centered on modeling, inference, and empirical behavior in the full routing setting, rather than as a paper whose main contribution is a regret theorem for a simplified variant.
        \item That said, our intention was not to replace all theory with experiments. The current paper already proves structural properties that are specific to the proposed formulation, including cross-expert information transfer and pruning invariance.
    \end{itemize}

    \item \textbf{Weakness 2 / Question 1: computational cost not clearly analyzed; runtime relative to baselines.}
    \begin{itemize}
        \item We agree that this needed to be made much clearer. In the fixed-parameter online-routing regime used on Jena and Melbourne, the dominant online cost of the proposed method can be summarized as
        \[
        T_t^{\mathrm{online}}
        =
        \widetilde O\!\left(
        C_\theta(x_t)
        + M d_g^3
        + M |\mathcal K_t| d_\alpha^3
        + |\mathcal E_t|\, C_{\mathrm{score}}(M,d_g,d_\alpha,S)
        \right),
        \]
        with online memory
        \[
        \mathrm{Mem}_t^{\mathrm{online}}
        =
        O\!\left(
        M d_g^2 + M |\mathcal K_t| d_\alpha^2
        \right).
        \]
        Here \(M\) is the number of regimes, \(d_g\) is the shared latent dimension, \(d_\alpha\) is the idiosyncratic state dimension, \(\mathcal E_t\) is the currently available set, \(\mathcal K_t\) is the maintained registry, and \(S\) is the Monte Carlo budget used in the \(g_z\) exploration rule. The important architectural point is that the expert-specific filtering and memory costs scale with the \emph{maintained registry} \(|\mathcal K_t|\), not with the cumulative catalog of all experts ever seen. Under pruning, \(|\mathcal K_t| \le |\mathcal E_t| + \Delta_{\max}\), so this part of the online state remains controlled when stale experts are removed.
        \item We also now report empirical runtime on the real-data benchmarks and completed runtime-oriented sweeps for the proposed method on both Jena and Melbourne. Table~\ref{tab:qba_runtime} shows that a substantial fraction of the runtime can be recovered by choosing better operating points, while preserving the best observed Jena performance and slightly improving Melbourne performance.
    \end{itemize}

    \begin{table}[h]
    \centering
    \caption{Completed operating-point sweeps for the proposed method. Lower is better for mean cost and runtime.}
    \label{tab:qba_runtime}
    \begin{tabular}{lccc}
    \toprule
    \textbf{Dataset / config} & \textbf{Mean cost} & \textbf{Runtime (ms/step)} & \textbf{Comment} \\
    \midrule
    Jena, original setting & 3.6100 & 3.225 & \(M=4, d_g=2, g_z, 25\) MC \\
    Jena, tuned setting & 3.6100 & 1.479 & \(M=4, d_g=1, g\) \\
    Melbourne, original setting & 5.7093 & 3.223 & \(M=4, d_g=2, g_z, 25\) MC \\
    Melbourne, tuned setting & 5.6959 & 1.688 & \(M=2, d_g=1, g_z, 12\) MC \\
    \bottomrule
    \end{tabular}
    \end{table}

    \item \textbf{Weakness 3 / Question 2: limited baseline comparison; why no Thompson / ensemble / shared-parameter LinUCB?}
    \begin{itemize}
        \item We thank the reviewer for this suggestion; we have directly addressed it. We added three stronger baselines to the full pipeline:
        \begin{itemize}
            \item Shared-Parameter LinUCB,
            \item Linear Thompson Sampling,
            \item Ensemble Sampling.
        \end{itemize}
        We then reran the real-data comparisons with this expanded baseline set. The resulting comparisons are summarized in Table~\ref{tab:qba_perf}. The empirical comparison is therefore no longer limited to LinUCB and NeuralUCB, and the proposed method remains the best adaptive method on both completed real-data benchmarks.
        \item We also agree that the shared-parameter LinUCB comparison is important. The original per-expert LinUCB was intended as a heterogeneous-expert baseline. In response to the review, we added a shared-parameter linear comparator using one global parameter vector over expert-conditioned features. This is the relevant common-parameter analogue in the present routing problem: a literal model with the same context feature vector for every expert would assign identical predictions to all experts and would therefore not define a meaningful routing baseline. The new shared-parameter LinUCB is stronger than the original per-expert LinUCB on the real-data benchmarks, and the proposed method still outperforms it.
    \end{itemize}

    \begin{table}[h]
    \centering
    \caption{Completed real-data comparisons with the expanded baseline set. Lower is better. We report mean cost; on Melbourne the completed run is over five seeds and on Jena the stochastic methods were also evaluated over five seeds.}
    \label{tab:qba_perf}
    \begin{tabular}{lcc}
    \toprule
    \textbf{Method} & \textbf{Jena} & \textbf{Melbourne} \\
    \midrule
    L2D-SLDS & 3.6100 & 5.6959 \\
    L2D-SLDS w/o \(g_t\) & 3.6633 & 5.7540 \\
    Shared-Parameter LinUCB & 3.9220 & 5.9582 \\
    Linear Thompson Sampling & 4.1197 & 5.9097 \\
    Ensemble Sampling & 4.0977 & 5.9290 \\
    NeuralUCB & 3.9755 & 5.9279 \\
    LinUCB & 5.1433 & 6.0998 \\
    \bottomrule
    \end{tabular}
    \end{table}
\end{itemize}


\section*{Official Review by Reviewer iNRb}

\textbf{Summary.}

The paper proposes a factorized switching linear dynamical system model for the learning-to-defer problem with dependent data. The problem seeks to randomize among experts to minimize cost in predictions where the choice yields a bandit feedback in that only the cost of the expert chosen is revealed post that choice. The expert's prediction costs are dependent across experts and time.

The idea of the construction of a dynamical system model for belief construction for the experts' prediction cost is well conceived. These costs are correlated across each other and over time. Modeling it as a switching linear dynamical system allows the authors to address the learning-to-defer problem with correlated data. The linear-Gaussian setup has the advantage of predictions to be all based on Gaussian updates, which is a plus.

There are perhaps four features of the whole model:
\begin{enumerate}
    \item Context-dependent regime switching
    \item Linear dynamics with two types of factors (global and expert-dependent)
    \item Information-directed sampling for exploration of experts
    \item Dynamic registry management
\end{enumerate}

The whole setup covers all bases of various modeling needs that this problem domain might need. Since the paper presents a new method without theoretical analysis, a more comprehensive experimental validation would have made a better case.

In practice, the predicted costs require model parameters such as , etc., and these would need to be estimated from data. The main body of the paper does not talk about estimation at all. It is presented in Algorithm 3 in the appendix. Once you factor in an estimation phase, the sample complexity will depend crucially on the size of the model. With that in mind, the whole idea of regimes and its soft-max-based model appears somewhat of an overkill. A discussion to that effect should be included.

The paper spends a lot of real estate in ``developing'' the model, but it does not quite provide a straightforward set of algorithms to implement it. It would have been easier if the linear-Gaussian model was presented way before and then it was explained.

\textbf{Soundness:} 3: good

\textbf{Presentation:} 2: fair

\textbf{Significance:} 3: good

\textbf{Originality:} 2: fair

\textbf{Key Questions For Author.}

Please see my comments about the strengths and weaknesses regarding the model complexity.

\textbf{Limitations:} No

\textbf{Rating:} 3: Weak reject: A paper with clear merits, but also some weaknesses, which overall outweigh the merits. Papers in this category require revisions before they can be meaningfully built upon by others. Please use sparingly.

\textbf{Confidence:} 3: You are fairly confident in your assessment. It is possible that you did not understand some parts of the submission or that you are unfamiliar with some pieces of related work. Math/other details were not carefully checked.

\subsection{Response}
\begin{itemize}
    \item The whole setup covers all bases of various modeling needs that this problem domain might need. Since the paper presents a new method without theoretical analysis, a more comprehensive experimental validation would have made a better case.
    \begin{itemize}
        \item As explained to reviewer @QbaF, deriving regret bounds is not possible because of the complexity of
        the setting we are working on. We however, give some theoretical implication from our model. As you
        suggested we have new experiments to give more depth to the paper: one on a real dataset Jane given in the
        general message and another one to explain that our method scales well when the number of experts explodes
        over time.
    \end{itemize}
    \item In practice, the predicted costs require model parameters such as , etc., and these would need to be
    estimated from data. The main body of the paper does not talk about estimation at all. It is presented in Algorithm 3 in the appendix. Once you factor in an estimation phase, the sample complexity will depend crucially on the size of the model. With that in mind, the whole idea of regimes and its soft-max-based model appears somewhat of an overkill. A discussion to that effect should be included.
    \begin{itemize}
        \item We totally agree that our approach relies on correctly estimating parameters like we believe any
        switching space type of approaches. Since this was not the heart of our approach (the estimation problem can
        be done with any estimation approach), we did not insist on it in the main aper. However, as you said, we
        give the algorithm in appendix. We will make sure to make it clearer in the final version of the paper.
    \end{itemize}
    \item The paper spends a lot of real estate in ``developing'' the model, but it does not quite provide a straightforward set of algorithms to implement it. It would have been easier if the linear-Gaussian model was presented way before, and then it was explained.
    \begin{itemize}
        \item We do believe that we have given the full detailed line by line algorithm defined in 'Algorithm 1
        Context-Aware Router L2D-SLDS' with corresponding fully detailed subroutines in Algorithm 2 & 3. We will
        clarify this in the main paper.
        \item We agree that introducing the linear system at first would have probably been a bit easier to follow.
        We will change the ordering to make it more logical. We thank the reviewer for this good suggestion.
    \end{itemize}
\end{itemize}


\section*{Official Review by Reviewer W5VY}

\textbf{Summary.}

The paper studies the problem of learning to defer in non-stationary time series settings with partial feedback and time-varying expert availability. The goal is to determine, at each round and for each context, which expert to query so that the cumulative cost incurred by these queries is minimized over time. The authors propose modeling the experts' predictions using a linear Gaussian dynamic model that incorporates both a shared global factor across experts and individual expert-specific factors. This model is combined with an IDS-inspired routing rule to minimize cumulative cost. The authors also conduct experiments to demonstrate the effectiveness of their method.

\textbf{Strengths.}
\begin{enumerate}
    \item The authors propose using a linear Gaussian dynamic model to represent the experts' predictions. This time-series perspective for modeling experts in bandit settings has been relatively less explored.
    \item The authors further introduce a model that includes both a shared global factor to capture cross-expert dependencies and individual factors to account for heterogeneity across experts. This modeling choice is also relatively underexplored in the literature and contributes to the novelty of the proposed method.
    \item The proposed methodology achieves lower cumulative cost in the numerical experiments, demonstrating its potential effectiveness.
\end{enumerate}

\textbf{Weaknesses.}
\begin{enumerate}
    \item I have two main concerns regarding the paper. The first relates to presentation. I found the paper difficult to read and hard to follow. This stems from several reasons. First, the paper introduces many technical terms and abbreviations without sufficient explanation, which hinders accessibility for a broader audience. For example, the abstract includes abbreviations such as L2D-SLDS and IDS without spelling out their full names. Moreover, some key modeling components are introduced without clear motivation. For instance, the shared factor and individual states are introduced at the beginning of Section 4.2, but their connection to the problem only becomes clear much later in Equation (13).
    \item My second concern relates to the evaluation. As a methodology paper, the work would benefit from either stronger theoretical guarantees or more comprehensive empirical validation. While the paper provides some theoretical analysis, it does not establish learning-theoretic results such as cumulative regret bounds or minimax optimality. The empirical evaluation is also somewhat limited. The synthetic experiment is relatively simple, relying on a linear AR(1) model, and the comparisons are restricted to classical bandit approaches such as LinUCB and NeuralUCB. Overall, neither the theoretical results nor the empirical evaluation fully justify the proposed methodology.
    \item While $C_k(x, y)$ is explicitly defined in Equation (1), $C_I(x, y)$ is not clearly defined when $I$ is a set. It is unclear whether it represents the sum of the individual costs for $k \in I$, or a vector containing all such individual costs.
\end{enumerate}

\textbf{Soundness:} 3: good

\textbf{Presentation:} 2: fair

\textbf{Significance:} 2: fair

\textbf{Originality:} 3: good

\textbf{Key Questions For Author.}
\begin{enumerate}
    \item What is the definition of $C_I$ on Line 81, Page 2?
    \item Would it be possible to establish an information-theoretic lower bound for the cumulative cost under the proposed setting, and to determine whether the regret of the proposed method matches this lower bound?
\end{enumerate}

\textbf{Limitations.}

I might have missed something, but I did not find adequate discussion of the limitations of the proposed methodology in the paper.

\textbf{Rating:} 3: Weak reject: A paper with clear merits, but also some weaknesses, which overall outweigh the merits. Papers in this category require revisions before they can be meaningfully built upon by others. Please use sparingly.

\textbf{Confidence:} 2: You are willing to defend your assessment, but it is quite likely that you did not understand the central parts of the submission or that you are unfamiliar with some pieces of related work. Math/other details were not carefully checked.


\subsection{Response}
\begin{itemize}
    \item I have two main concerns regarding the paper. The first relates to presentation. I found the paper difficult to read and hard to follow. This stems from several reasons. First, the paper introduces many technical terms and abbreviations without sufficient explanation, which hinders accessibility for a broader audience. For example, the abstract includes abbreviations such as L2D-SLDS and IDS without spelling out their full names. Moreover, some key modeling components are introduced without clear motivation. For instance, the shared factor and individual states are introduced at the beginning of Section 4.2, but their connection to the problem only becomes clear much later in Equation (13).
    \begin{itemize}
        \item We thank the reviewer for its suggestion. We will make sure that such abbreviations will be explained
        more in details to be more welcomed for a broader audience. We will also work on restructuring the
        modelisation part; specifically starting with the model formulation and then explain term by term what each
        terms is doing and what is the intuition behind.
    \end{itemize}
\end{itemize}





